\documentclass[11pt]{article}
\usepackage[a4paper,margin=2.3cm]{geometry}
\usepackage{amsmath,amssymb}

\usepackage{fontspec}
\usepackage{booktabs,array}
\usepackage[hidelinks]{hyperref}
\linespread{1.12}
\setlength{\parskip}{0.45em}
\setlength{\parindent}{0pt}
\newcommand{\cl}{c_{\text{late}}}
\newcommand{\cs}{c_{\text{slack}}}
\newcommand{\1}{\mathbf 1}
\newcounter{prop}
\newenvironment{prop}{\refstepcounter{prop}\par\medskip\noindent\textbf{Proposition \theprop.}\ \itshape}{\par\medskip}
\newenvironment{pf}{\par\noindent\textbf{Proof.}\ }{\par\noindent\hfill QED\par\medskip}
\begin{document}

{\Large\textbf{Why a quantile sets the promised date, and where the calibration guarantee stops}}

\medskip
This note answers two questions. First, why the promised date is a quantile of delivery time rather than its mean. Second, what ``the coverage is met'' means after calibration. I separate three quantities that are easy to confuse: a proportion on the calibration set, a probability under exchangeability, and a proportion observed in the evaluation period. The results used here are standard: the quantile as the minimiser of an asymmetric linear loss, and the rank argument behind split conformal calibration. I state and prove them in the form this project uses and then apply them; none of them is new.

\bigskip
\textbf{1. Why an asymmetric loss leads to a quantile}

Let the delivery time be an integer \(Y\ge1\) with \(\mathbb E[Y]<\infty\), and the promise an integer \(d\ge1\). A promise that is too short makes the order late; one that is too long leaves the customer waiting. I describe both costs with a linear loss,
\[
L(d,Y)=\cl\,(Y-d)_+ +\cs\,(d-Y)_+,\qquad \cl,\cs>0 .
\]
This is a simplifying assumption: cost grows linearly in days. Let \(R(d)=\mathbb E[L(d,Y)]\), \(F(d)=P(Y\le d)\) and \(q=\cl/(\cl+\cs)\).

\begin{prop}
\(d^\ast=\min\{d\ge1:\ F(d)\ge q\}\) minimises \(R(d)\). If \(F(d^\ast)=q\), then \(d^\ast+1\) is also optimal.
\end{prop}

\begin{pf}
Compare a promise of \(d\) days with one of \(d+1\). Order by order,
\[
L(d+1,Y)-L(d,Y)=
\begin{cases}
-\cl, & Y>d,\\
+\cs, & Y\le d .
\end{cases}
\]
Taking expectations,
\[
\Delta(d):=R(d+1)-R(d)=\cs F(d)-\cl\bigl(1-F(d)\bigr)=(\cl+\cs)\bigl(F(d)-q\bigr).
\]
\(F\) is non-decreasing, so \(\Delta\) is non-decreasing. For \(d<d^\ast\), \(F(d)<q\), so \(\Delta(d)<0\); for \(d\ge d^\ast\), \(\Delta(d)\ge0\). Thus \(R\) strictly decreases up to \(d^\ast\) and does not decrease afterwards, so \(d^\ast\) is a minimiser. Since \(F(0)=0<q\), \(d^\ast\ge1\). If \(F(d^\ast)=q\), then \(\Delta(d^\ast)=0\), that is, \(R(d^\ast+1)=R(d^\ast)\).
\end{pf}

For a single order with information \(X=x\) at checkout, replace \(F\) by \(F(d\mid x)=P(Y\le d\mid X=x)\); the optimal promise is then the conditional \(q\)-quantile. With \(\cl=q\) and \(\cs=1-q\), \(L\) is the quantile (pinball) loss.

This loss gives the quantile target its decision meaning, but it is not every model's fitting objective. In the project the linear model fits a quantile regression on the \(\log(1+Y)\) scale and transforms back to days; HGB fits on days directly; model selection and evaluation are on the day scale. Quantiles are preserved under monotone transformations, so both scales aim at the same conditional quantile. In a finite sample, however, the two fitting objectives differ.

The mean minimises squared loss, which treats a late day and a spare day as equally costly. The mean does not generally minimise the asymmetric loss defined here; for some distributions it may happen to coincide with the optimal quantile. Choosing \(q=0.90\) is the same as assuming \(\cl/\cs=9\): one late day costs as much as nine spare days. This is a target I chose, not a cost ratio estimated from the Olist data.

\bigskip
\textbf{2. Calibration: three different ``coverages''}

A model's \(f(x)\) is only an estimate, so using it directly need not reach the target. A separate calibration set corrects it. Let \(f\) be fitted already, without the calibration set. On \(n\) calibration orders compute
\[
r_i=Y_i-f(X_i),\qquad k=\lceil (n+1)q\rceil\le n,\qquad s=r_{(k)},\qquad U(x)=f(x)+s ,
\]
where \(r_{(k)}\) is the \(k\)-th smallest residual. Note that \(Y_i\le U(X_i)\iff r_i\le s\).

\textbf{First quantity: the proportion on the calibration set, at least \(k/n\).}

\begin{prop}
For any data, \(\displaystyle \hat c_{\mathrm{cal}}:=\frac1n\sum_{i=1}^n \1\{Y_i\le U(X_i)\}\ \ge\ \frac kn\).
\end{prop}

\begin{pf}
By definition, at least \(k\) residuals are at most \(r_{(k)}=s\) (more if there are ties). The claim follows from \(Y_i\le U(X_i)\iff r_i\le s\).
\end{pf}

This inequality uses no probability assumption, and it holds even if delivery conditions later change. It is a necessary consequence of the specified calibration rule. Satisfying it alone neither verifies the implementation nor establishes coverage on future orders: a shift wrongly set to 1,000 days would satisfy it too. A more direct implementation check is to recompute the \(k\)-th residual independently and compare it with the frozen shift.

\textbf{Second quantity: the probability that a new order is covered, at least \(k/(n+1)\).}

\begin{prop}
If \((X_1,Y_1),\dots,(X_n,Y_n),(X_{n+1},Y_{n+1})\) are exchangeable, then
\[
P\bigl(Y_{n+1}\le U(X_{n+1})\bigr)\ \ge\ \frac{k}{n+1}\ \ge\ q .
\]
\end{prop}

\begin{pf}
Since \(f\) is fixed, \(r_1,\dots,r_{n+1}\) are exchangeable too. Break ties at random and rank the \(n+1\) residuals; let \(\rho\) be the rank of the new residual \(r_{n+1}\). By exchangeability \(\rho\) is uniform on \(\{1,\dots,n+1\}\), so \(P(\rho\le k)=k/(n+1)\). If \(\rho\le k\), at most \(k-1\) of the first \(n\) residuals rank before \(r_{n+1}\); hence the \(k\)-th smallest of them is at least \(r_{n+1}\), that is, \(r_{n+1}\le s\). Thus \(P(r_{n+1}\le s)\ge k/(n+1)\). Finally, \(k\ge(n+1)q\).
\end{pf}

The denominator is \(n+1\) because the new order is itself one of the \(n+1\) exchangeable draws. The probability averages over which calibration set was drawn and which new order arrives. It does not guarantee the target in every month or state, nor after one particular calibration set has been observed. Random tie-breaking is used only in the proof; the algorithm is unchanged.

\textbf{Third quantity: the proportion observed in the evaluation period.}

With \(m\) evaluation orders,
\[
\hat c_{\mathrm{eval}}:=\frac1m\sum_{j=1}^m \1\{Y_j\le U(X_j)\}.
\]
This is an observed proportion, not a probability. If every evaluation order is exchangeable with the calibration orders, Proposition 3 gives \(\mathbb E[\hat c_{\mathrm{eval}}]\ge k/(n+1)\); even then, a finite \(m\) adds sampling variation, so a single value can fall below or above it. If exchangeability fails, it describes only these evaluation orders.

\begin{center}
\small
\begin{tabular}{>{\raggedright\arraybackslash}p{0.25\textwidth}>{\raggedright\arraybackslash}p{0.33\textwidth}>{\raggedright\arraybackslash}p{0.33\textwidth}}
\toprule
Quantity & What it shows & What it does not show\\
\midrule
\(\hat c_{\mathrm{cal}}\ge k/n\) (Prop.~2) & A necessary consequence of the specified rule & That the implementation is correct; that future orders will reach the target\\
\(P(Y_{n+1}\le U)\ge k/(n+1)\) (Prop.~3) & Under exchangeability, coverage of at least \(q\), averaged over the calibration set and the new order & The target in every calibration set, state or month. Exchangeability between C and E has not been established, so the guarantee is not asserted for the evaluation period\\
\(\hat c_{\mathrm{eval}}\) & Empirical performance of the frozen policy on these evaluation orders & Reliability on all orders or in future operations\\
\bottomrule
\end{tabular}
\end{center}

A simple example shows that the first and third quantities can come apart completely. The model always predicts 10 days, and every calibration order arrives in 10 days, so \(s=0\) and \(\hat c_{\mathrm{cal}}=100\%\). If future orders all take 15 days, \(\hat c_{\mathrm{eval}}=0\). Conversely, if they all take 5 days, \(\hat c_{\mathrm{eval}}=100\%\), but every order is promised 5 days too many. Proposition 2 holds in both cases; the exchangeability behind Proposition 3 does not, so its guarantee does not carry over (\href{https://arxiv.org/abs/2202.13415}{Barber et al.}).

\bigskip
\textbf{3. Whole-day promises}

A promise must be a whole number of days. Since \(Y\) is an integer, \(Y\le U\iff Y\le\lfloor U\rfloor\). Thus rounding down covers exactly the orders that \(U\) covers; a minimum of one day can only enlarge that set. The registered rule is \(\max\{1,\lceil U\rceil\}\). It covers every order that \(U\) covers, so the lower bounds of Propositions 2 and 3 hold for it as well. Rounding \(U\) up to \(\lceil U\rceil\) adds less than one day. The minimum-one-day constraint is a separate adjustment and can only increase coverage.

\bigskip
\textbf{4. Back to Olist}

The calibration set has \(n=17{,}056\) orders and \(k=15{,}352\), so \(k/n\approx90.01\%\) and \(k/(n+1)\approx90.00\%\). For the linear model the three quantities are:
\begin{itemize}\setlength{\itemsep}{0pt}
\item on the calibration set, the registered rounded-up promise covers 91.19\%, not below \(k/n\), consistent with Proposition 2; recomputing the \(k\)-th residual independently gives the frozen shift (6.1479 days);
\item the guarantee under exchangeability is at least \(k/(n+1)\), but whether C (October--December 2017) and the evaluation period E (April--July 2018) are exchangeable is itself the question;
\item in the evaluation period, the observed coverage on 19,950 orders is 96.18\%.
\end{itemize}

The over-coverage on E suggests that the relation between the calibration and evaluation periods deserves a closer look. It is compatible with changes in delivery times, with the sample selection, and with changes in the order mix. This note runs no formal test of exchangeability and cannot separate these contributions. Therefore I treat it as a diagnostic signal, not as confirmation of any single cause. All conclusions apply only to single-seller orders that were eventually delivered.

\end{document}
